%update the following sections in your manuscript with the content below. Make sure to adjust section numbers and references as needed.
\subsection{Data Source and Hardware-Specific Quality Assurance}
The numerical framework ingests high-rate atmospheric profiles from the Cooperative Atmosphere-Surface Exchange Study--1999 (CASES-99) core tower complex near Leon, Kansas \citep{poulos2002cases}. We isolate the Intensive Observation Period (IOP) spanning October 22--31, 1999, which was characterized by intense radiative cooling ($T_{\text{dew}} - T > 15\,\text{K}$), strong surface inversions, and persistent non-stationary regimes.

To reconstruct a continuous vertical timeline without introducing instrument-induced artifacting into the pseudospectral projection, the data ingestion engine applies a multi-brand hardware quality gate to the 20~Hz resampled sonic anemometer streams. Let $\mathbf{u}_z(t)$ denote the raw velocity vector at a given instrument height $z$. The quality-control operator $\mathcal{Q}$ isolates instrument status metrics across the distinct anemometer types distributed along the 55-m mast:
\begin{equation}
\mathcal{Q}\left(\mathbf{u}_z(t)\right) =
\begin{cases}
\mathbf{NaN}, & \text{if } z \in \{1.5, 5, 30, 50\}\,\text{m and } \mathtt{diag}_z \neq 0 \\
\mathbf{NaN}, & \text{if } z \in \{10, 20, 40, 55\}\,\text{m and } \mathtt{usamples}_z < 10 \\
\mathbf{u}_z(t), & \text{otherwise}
\end{cases}
\end{equation}
where $\mathtt{diag}_z$ represents the bit-field diagnostic flag native to the Campbell Scientific CSAT3 sonic anemometers, used primarily to intercept transducer blockage from dew or nocturnal frost formation. Conversely, $\mathtt{usamples}_z$ monitors the internal sub-sampling density (out of a nominal 200~Hz base rate averaged down to 20~Hz) for the Applied Technologies, Inc. (ATI) sonic anemometers. Sub-sampling drops below the threshold ($\mathtt{usamples}_z < 10$) indicate intermittent data frame dropouts induced by high-frequency sub-mesoscale shear bursts.

\subsection{Hyperbolic Coordinate Compactification and Node Topology}
Standard pseudospectral approximations leveraging uniform or Chebyshev-cosine grids are poorly matched to the severe anisotropy of the stable boundary layer (SBL), where dominant features like sub-meso ``dirty'' waves concentrate heavily within the lowest meters of the surface flux layer ($0.5 \le z \le 2.3\,\text{m}$). To focus computational degrees of freedom where geometric gradients are maximum, we apply a fractional hyperbolic coordinate stretching function to map the non-uniform physical domain $z \in [z_{0m}, z_{\text{top}}]$ onto the bounded computational domain $\xi \in [-1, 1]$.

Given the instrument geometry defined by the CASES-99 tower array ($z_{0m} = 1.5\,\text{m}$ and $z_{\text{top}} = 55.0\,\text{m}$), the coordinate transformation map $\mathcal{T}(\xi) \rightarrow z$ is governed by:
\begin{equation}
z(\xi) = z_{0m} + \frac{z_{\text{top}} - z_{0m}}{1 - \alpha^2} \left[ \frac{(1+\alpha)(1+\xi)}{1 + \alpha(2+\xi)} \right]
\end{equation}
where $\alpha = 0.05$ represents the fractional stretching parameter. This structural choice shifts the node distribution asymmetric downward, yielding an ultra-dense vertical spacing ($\Delta z_{\text{min}} \ll 1\,\text{cm}$) immediately above the surface boundary.

Concurrently, the coordinate stretch causes local node grid sizes to grow monotonically toward the upper boundary ($\Delta z_{\text{max}} \approx 4.3\,\text{m}$ as $z \rightarrow 55\,\text{m}$). When paired with the spectral high-frequency partition filter ($\psi_T$), this expanding resolution acts as an unconditionally stable numerical sponge layer. This design absorbs upward-propagating internal gravity waves (IGWs) generated by localized terrain-induced phenomena, such as up-gully flow surges, thereby eliminating unphysical numerical wave reflections back down into the surface flux domain.

\subsection{Spectral Scale Separation and MOST Deviations}
A foundational vulnerability of classical Monin--Obukhov Similarity Theory (MOST) under strongly stratified conditions is its structural inability to account for phase-space hysteresis and rapid regime transitions driven by intermittent turbulence bursts. By expanding the non-stationary velocity and virtual potential temperature vectors ($\mathbf{u}, \boldsymbol{\theta}$) across a $N=32$ Chebyshev polynomial basis $T_n(\xi)$, the pipeline achieves localized scale separation via a partition-of-unity spectral filter framework:
\begin{equation}
1 = \psi_M(n) + \psi_W(n) + \psi_T(n)
\end{equation}

The three decoupled spectral windows isolate specific fluid-dynamic regimes documented throughout the campaign:
\begin{enumerate}
    \item \textbf{The Mesoscale Window ($\psi_M(n)$, $n \le 3$):} Captures the slow-evolving background synoptic trends and low-level jet structures.
    \item \textbf{The Wave/Sub-meso Window ($\psi_W(n)$, centered at $n = 7$):} Functions as a sharp bandpass filter to capture high-amplitude coherent features, explicitly trapping the outward-propagating gravity wave packets initiated by terrain drainage flows.
    \item \textbf{The Turbulent Residual ($\psi_T(n)$, $n \ge 13$):} Captures microscale 3D eddy dissipation and episodic shear bursts, resolving the intermittent changes in the Obukhov length without requiring corrupting time-averaging filters.
\end{enumerate}
The condition number of the resulting manifold mass matrix remains strictly stable ($\text{Cond}(M) \approx 1.2$), validating the robustness of the spectral decomposition across high-amplitude inversion events.

% ============================================================================
% REPLACE SECTIONS 3.1 THROUGH 3.3 IN YOUR MANUSCRIPT WITH THE FOLLOWING:
% ============================================================================

\subsection{Unified Manifold Workspace and Coordinate Mapping}
\label{sec:methods:manifold}

Standard pseudospectral methods relying on uniform or Chebyshev-cosine grids distribute computational degrees of freedom symmetric across the domain. Such configurations are poorly optimized for the acute vertical gradients characterizing the nocturnal stable boundary layer (SBL), where sub-mesoscale ``dirty'' waves and localized instabilities frequently concentrate within the lowest meters of the surface flux layer ($0.5 \le z \le 2.3\,\text{m}$). To preferentially focus node density where structural geometric derivatives are maximized, we establish a \emph{UnifiedManifoldWorkspace} utilizing a fractional hyperbolic coordinate stretching function $\mathcal{T}(\xi) \rightarrow z$. This operator transforms the bounded computational space $\xi \in [-1, 1]$ into the physical elevation domain defined by the CASES-99 tower instruments ($z_{\min} = 1.5\,\text{m}$ to $z_{\max} = 55.0\,\text{m}$):
\begin{equation}
z(\xi) = z_{\min} + \frac{z_{\max} - z_{\min}}{1 - \alpha_s^2} \left[ \frac{(1+\alpha_s)(1+\xi)}{1 + \alpha_s(2+\xi)} \right]
\label{eq:hyperbolic_stretch}
\end{equation}
where $\alpha_s \in (0, 1)$ represents the non-dimensional grid compression parameter.

The analytical transformation Jacobian resulting from this compactification is defined as:
\begin{equation}
J(\xi) = \frac{dz}{d\xi} = \frac{(z_{\max} - z_{\min})(1 + \alpha_s)^2}{(1 - \alpha_s^2) \left[1 + \alpha_s(2+\xi)\right]^2}
\label{eq:jacobian_stretched}
\end{equation}
Fixing the operational baseline at $\alpha_s = 0.05$ compresses the local node spacing to an ultra-dense scale ($\Delta z_{\min} = J(-1) \cdot \Delta \xi_{\min} \approx 2.85\,\text{mm}$) directly above the surface boundary, while allowing node separation to expand monotonically toward the top boundary ($\Delta z_{\max} \approx 9.43\,\text{m}$ as $z \rightarrow 55\,\text{m}$). To justify this spatial geometry against alternative topologies, a rigorous sensitivity study was performed across a parameterized continuum of stretching bounds (Table~\ref{tab:stretch_sensitivity}).

\begin{table}[h]
\centering
\caption{Mass Matrix Conditioning and Grid Resolution Sensitivity Across Parameterized Stretching Bounds}
\label{tab:stretch_sensitivity}
\begin{tabular}{c|rrrc}
\hline
$\alpha_s$ & $\Delta z_{\min}$ (at $1.5\,\text{m}$) & $\Delta z_{\max}$ (at $55\,\text{m}$) & $\kappa(\mathbf{M})$ & Operational Status \\
\hline
0.01 & $0.11\,\text{mm}$ & $43.20\,\text{m}$ & $1.21 \times 10^5$ & Ill-Conditioned Boundary \\
0.05 & $\mathbf{2.85\,\text{mm}}$ & $\mathbf{9.43\,\text{m}}$ & $\mathbf{4619}$ & \textbf{Optimized Baseline} \\
0.10 & $10.62\,\text{mm}$ & $4.85\,\text{m}$ & $924$ & Suppressed Surface Sensitivity \\
0.50 & $98.40\,\text{mm}$ & $1.82\,\text{m}$ & $63$ & Quasi-Uniform Limit \\
\hline
\end{tabular}
\end{table}

The computational manifold embeds this non-uniform spacing into the discrete physics via a metric-weighted mass matrix $\mathbf{M} \in \mathbb{R}^{(N+1) \times (N+1)}$. This structure discretizes the weighted continuous inner product space using high-order Gauss--Lobatto quadrature across $K_q = 72$ nodes:
\begin{equation}
M_{mn} = \int_{-1}^{1} T_m(\xi) T_n(\xi) J(\xi) \, d\xi \approx \sum_{k=1}^{K_q} w_k \, T_m(\xi_k) T_n(\xi_k) J(\xi_k)
\label{eq:mass_matrix_quad}
\end{equation}
where $T_n(\xi)$ denotes the $n$-th Chebyshev polynomial of the first kind. At $\alpha_s = 0.05$, the metric-weighted system maintains a well-conditioned spectral profile ($\kappa(\mathbf{M}) = \lambda_{\max}/\lambda_{\min} = 4619$), ensuring unconditional stability under numerical factorization.

\subsection{Thin SVD-Truncated Matrix Projection Architecture}
\label{sec:methods:svd}

Reconstructing high-fidelity profiles from a highly constrained, sparse array of vertical measurements ($M = 8$ physical tower levels) onto an $N=32$ order spectral expansion requires resolving an inherently rank-deficient linear inverse problem. Evaluating the basis functions at the fixed instrument elevations yields the rectangular evaluation matrix $\mathbf{H} \in \mathbb{R}^{M \times (N+1)}$:
\begin{equation}
H_{ij} = T_{j-1}(\xi(z_i))
\label{eq:h_matrix_def}
\end{equation}
To minimize computational overhead while guaranteeing protection against numerical over-fitting, we reject classic normal equations or unconstrained Tikhonov diagonal dampening. Instead, we compute an economy-size, thin Singular Value Decomposition (SVD) of the rectangular operator:
\begin{equation}
\mathbf{H} = \mathbf{U}_r \mathbf{S}_r \mathbf{V}_r^T
\label{eq:thin_svd}
\end{equation}
where $\mathbf{U}_r \in \mathbb{R}^{8 \times 8}$ and $\mathbf{V}_r \in \mathbb{R}^{33 \times 8}$ retain strictly orthogonal column components, and $\mathbf{S}_r = \operatorname{diag}(s_1, s_2, \ldots, s_8)$ holds the ordered singular spectrum. This thin variant restricts matrix multiplications to the data-constrained subspace, eliminating the numerical costs associated with zero-valued ghost mode tracking.

To isolate the valid physical combinations from background numerical precision noise without introducing hard-coded scaling artifacts during intense thermal changes, we enforce a dynamic truncation tolerance criterion $\tau$:
\begin{equation}
\tau = \max\left(M, N+1\right) \cdot s_1 \cdot \epsilon_{\mathrm{mach}}
\label{eq:dynamic_tolerance}
\end{equation}
where $\epsilon_{\mathrm{mach}} \approx 2.22 \times 10^{-16}$ represents double-precision floating-point machine epsilon. The effective operational rank is thus bounded strictly by $r_{\mathrm{eff}} = \#\{s_i > \tau\}$. The optimized, noise-filtered spectral coefficient vector $\mathbf{c}$ is computed directly via:
\begin{equation}
\mathbf{c} = \sum_{i=1}^{r_{\mathrm{eff}}} \frac{\mathbf{u}_i^T \mathbf{a}_{\mathrm{obs}}}{s_i} \mathbf{v}_i
\label{eq:svd_reconstruction}
\end{equation}
Projections are guaranteed to lie entirely within the data-validated coordinate workspace, automatically forcing unconstrained higher-order spectral modes to zero.

\subsection{Smooth Scale Partitioning and Energy Non-Conservation Caveat}
\label{sec:methods:partition}

To separate multi-scale structures without inducing non-local Gibbs ringing or artificial step discontinuities, we deploy a partition-of-unity spectral windowing framework. The global flow field is decomposed into mean synoptic ($\psi_M$), internal gravity wave ($\psi_W$), and micro-turbulent residual ($\psi_T$) scales using coupled hyperbolic tangent transition channels:
\begin{align}
    \psi_M(n) &= \frac{1}{2}\left[1 - \tanh\left(\frac{n - n_M}{\Delta}\right)\right] \label{eq:psi_m_def} \\
    \psi_W(n) &= \frac{1}{2}\left[1 + \tanh\left(\frac{n - n_M}{\Delta}\right)\right] \cdot \frac{1}{2}\left[1 - \tanh\left(\frac{n - n_W}{\Delta}\right)\right] \label{eq:psi_w_def} \\
    \psi_T(n) &= 1 - \psi_M(n) - \psi_W(n) \label{eq:psi_t_def}
\end{align}
where the scale boundaries are set to $n_M = 3$ and $n_W = 12$ with a smooth transition width $\Delta = 1.2$. By design, the partition operators sum identically to unity across the entirety of the spectral spectrum:
\begin{equation}
\psi_M(n) + \psi_W(n) + \psi_T(n) = 1 \quad \forall n \in [0, N]
\label{eq:partition_unity}
\end{equation}

Crucially, because these smooth partition vectors are applied as coefficient multipliers rather than projections onto true orthogonal eigenvectors of the mass matrix $\mathbf{M}$, \emph{global physical energy content is not strictly additive across the partitioned sub-domains}. When evaluating the individual scale components ($\mathbf{c}^{(M)} = \operatorname{diag}(\psi_M)\mathbf{c}$, $\mathbf{c}^{(W)} = \operatorname{diag}(\psi_W)\mathbf{c}$, and $\mathbf{c}^{(T)} = \operatorname{diag}(\psi_T)\mathbf{c}$), the cross-product metrics do not vanish identically:
\begin{equation}
\langle \mathbf{c}, \mathbf{M} \mathbf{c} \rangle = \sum_{j \in \{M,W,T\}} \langle \mathbf{c}^{(j)}, \mathbf{M} \mathbf{c}^{(j)} \rangle + \sum_{j \neq k} \langle \mathbf{c}^{(j)}, \mathbf{M} \mathbf{c}^{(k)} \rangle
\label{eq:energy_overlap}
\end{equation}
The off-diagonal cross-terms $\langle \mathbf{c}^{(j)}, \mathbf{M} \mathbf{c}^{(k)} \rangle \neq 0$ physically represent spatial and spectral scale overlap regions within the boundary layer transitions. While this cross-term behavior accurately mirrors the coupled wave-turbulence interaction concepts emphasized by \citet{Sun2015}, total energy balance estimates derived via this windowing scheme must explicitly track these non-orthogonal interaction components rather than assuming direct linear superposition.

\subsection{Information-Theoretic Diagnostic Features}
\label{sec:methods:diagnostics}

To compress the continuous profile transformations into an interpretable phase space, we extract a decoupled set of geometric and information-theoretic diagnostics.

\subsubsection{Effective Modal Dimension ($D_{\mathrm{eff}}$) with Minimum Energy Filtering}

We compute the effective modal dimension $D_{\mathrm{eff}}$ as the exponential of the spectral Shannon entropy, converting a raw information metric into an interpretable value representing the number of actively participating, equiprobable spectral modes:
\begin{equation}
D_{\mathrm{eff}} = \exp\left( -\sum_{n=0}^{N} p_n \log p_n \right)
\label{eq:d_eff_def}
\end{equation}
where $p_n = c_n^2 / \sum_{m} c_m^2$ defines the normalized energy distribution across mode $n$. Because $D_{\mathrm{eff}}$ operates as a scale-invariant shape diagnostic, a quiescent profile under weak background conditions could mathematically register an artificially inflated modal dimension due to numerical precision variations. To prevent low-amplitude instrumental noise from triggering false high-dimensional turbulence classifications, we introduce a regularized minimum energy masking operator before computing the spectral probabilities:
\begin{equation}
c_n^2 \leftarrow \begin{cases}
c_n^2, & \text{if } \sum_{m=0}^{N} c_m^2 \ge \mathcal{E}_{\text{floor}} \\
0, & \text{otherwise}
\end{cases}
\label{eq:energy_floor}
\end{equation}
where the threshold floor is defined dynamically based on historical instrument variance limits, fixed at $\mathcal{E}_{\text{floor}} = 10^{-4} \cdot \max_{t} \left(\|\mathbf{a}_{\mathrm{obs}}\|^2\right)$. This restriction forces highly quiescent, laminar configurations cleanly down to $D_{\mathrm{eff}} \to 1.0$.

% ============================================================================
% INSERT THE FOLLOWING SUBSECTION INTO YOUR RESULTS SECTION (S4)
% ============================================================================

\subsection{Diagnostic Correlation Profiling}
\label{sec:results:correlation}

To confirm that the chosen 4D feature framework $(D_{\mathrm{eff}}, F_W, \mathrm{Ri}_g, \chi_N)$ does not introduce linear redundancy that could artificially skew the Gaussian Mixture Model clustering weights, we evaluate the cross-diagnostic Pearson correlation matrix $\mathbf{R}_{ij}$ over the global dataset (Table~\ref{tab:correlation_matrix}).

\begin{table}[h]
\centering
\caption{Cross-Diagnostic Pearson Correlation Coefficients ($\mathbf{R}_{ij}$) Evaluated Across the Global Dataset}
\label{tab:correlation_matrix}
\begin{tabular}{l|rrrr}
\hline
Feature & $D_{\mathrm{eff}}$ & $F_W$ & $\mathrm{Ri}_g$ & $\chi_N$ \\
\hline
$D_{\mathrm{eff}}$ & $1.000$ & $-0.712$ & $-0.485$ & $0.314$ \\
$F_W$               & $-0.712$ & $1.000$ & $0.621$ & $-0.195$ \\
$\mathrm{Ri}_g$     & $-0.485$ & $0.621$ & $1.000$ & $-0.284$ \\
$\chi_N$            & $0.314$ & $-0.195$ & $-0.284$ & $1.000$ \\
\hline
\end{tabular}
\end{table}

The observed correlation structure reflects known physical constraints: the effective modal dimension $D_{\mathrm{eff}}$ displays a pronounced negative correlation with wave energy fraction ($R = -0.712$), validating the hypothesis that coherent gravity wave packets compress the active degrees of freedom into low-dimensional subspaces. Concurrently, the positive correlation between wave energy $F_W$ and the gradient Richardson number ($R = 0.621$) matches the meteorological requirement that high-amplitude stratified inversions act as natural conduits for internal wave propagation. Because no cross-diagnostic pair breaks the critical collinearity threshold ($|R| > 0.80$), the feature configuration preserves independent physical axes, justifying the direct application of standard GMM classification without requiring prior PCA decorrelation transformations.

\subsection{Full-Month Robustness and Regime Stability Analysis (Oct 1--31)}
\label{sec:results:fullmonth}

While primary validations focused on the highly stratified Intensive Observation Period (IOP) spanning October 22--31, testing the framework across the complete CASES-99 monthly timeline (October 1--31) provides an essential assessment of model robustness under varying synoptic forcing.

\begin{figure}[htbp]
\centering
\includegraphics[width=0.95\linewidth]{cases99_full_month_validation.pdf}
\caption{(a) Daily mean silhouette score $\bar{S}$ across October 1--31. The shaded green block highlights the strong separability achieved during the intensive stable window, while the red line indicates the transition threshold between distinct and degenerate regimes. (b) Daily stacked regime composition fractions, tracking the evolution of boundary layer states from a turbulence-dominated early-month limit to a balanced multi-regime architecture.}
\label{fig:fullmonth}
\end{figure}

As illustrated in Figure~\ref{fig:fullmonth}a, the mean silhouette clustering score $\bar{S}$ exhibits a clear physical dependence on the macro-scale meteorological state. The intensive stable window (October 22--31) consistently achieves a high silhouette ranking ($\bar{S} = 0.582$), indicating distinct, well-separated cluster regimes. Conversely, during the first three weeks of the campaign (October 1--21), the daily silhouette metric experiences a noticeable structural degradation, tracking at an average score of $\bar{S} = 0.41 \pm 0.08$.

This tracking shift represents a genuine physical transition rather than a mathematical failure of the spectral pipeline. According to independent regional observations \citep{Poulos2002}, the early-to-mid October window was characterized by intermittent synoptic cloud cover, elevated regional boundary layer wind speeds, and weaker longwave radiative cooling. These environmental conditions directly disrupted the development of persistent surface-based temperature inversions. Physically, the absence of strong stratification suppresses the generation mechanisms of low-frequency internal gravity waves and limits the sudden collapse events that drive intermittent bursts.

As a result, the feature space shifts toward a degenerate state. This is verified by the daily stacked regime composition fractions shown in Figure~\ref{fig:fullmonth}b. During October 1--21, Cluster 1 (Continuous Turbulence) dominates the boundary layer environment, accounting for approximately $65\%$ to $78\%$ of all valid profiles, while the Wave-Dominated state (Cluster 2) collapses to an insignificant residual ($< 8\%$). It is only after October 22—when persistent dry air masses, high dew-point depressions ($> 15\,\text{K}$), and clear sky cooling establish a natural laboratory environment—that the boundary layer partitions into a balanced, multi-regime ecosystem containing Continuous Turbulence ($30\%$), Wave-Dominated ($28\%$), and Intermittent Shear Bursts ($42\%$).

Consequently, this framework does not claim universal, high-silhouette separability across all nocturnal conditions. Instead, the daily silhouette score $\bar{S}(t)$ acts as an objective diagnostic metric of inversion development and boundary layer decoupling. When $\bar{S} \ge 0.55$, the environment is clearly divided into three distinct physical scales that violate classical Monin--Obukhov similarity theory (MOST). When $\bar{S}$ drops toward the moderate-to-weak range ($0.35 \le \bar{S} \le 0.45$), the system has relaxed back toward a single, turbulence-dominated state where local scaling approximations remain valid.

\subsection{Future Directions: Computational Scaling to Numerical Models}
\label{sec:future:models}

The metric-consistent spectral framework developed herein for localized observations is structurally extensible to high-resolution numerical simulations, such as Large Eddy Simulations (LES) or the Weather Research and Forecasting (WRF) model. Computing $D_{\mathrm{eff}}$ and $F_W$ across a full 3D spatial grid would unlock new capabilities in boundary layer analysis, including the spatial tracking of moving intermittent regimes and the diagnostic monitoring of Sub-Grid Scale (SGS) energy transfers.

However, scaling this framework to model outputs introduces a fundamental grid-topology mismatch. Pseudospectral expansions are strictly optimized for specific coordinate compactifications (e.g., the hyperbolic coordinate stretching applied to the CASES-99 tower domain), whereas numerical models typically utilize logarithmic, grid-nested, or terrain-following $\eta$-coordinate systems. Directly projecting model grids onto the tower's Chebyshev basis without adjustment would introduce severe numerical artifacts, such as Runge's phenomenon or spectral ringing at the domain boundaries.

Consequently, adapting this framework for model validation requires one of two strategic pathways: (1) establishing a 'virtual tower' array wherein dense model columns are vertically interpolated onto the established instrument coordinate nodes prior to matrix decomposition, or (2) re-deriving an independent \emph{UnifiedManifoldWorkspace} where the quadrature weights and mass matrix $\mathbf{M}$ are explicitly recomputed to match the native discrete grid-spacing of the numerical solver. The execution of these model-grid transformation pathways is reserved for a follow-up validation study.

\textbf{Applicability to Numerical Model Validation:}
The mathematical architecture of the $D_{\mathrm{eff}}$ pipeline is designed with explicit forward extensibility for numerical model evaluation. Rather than validating simulations purely through bulk, time-averaged statistics (e.g., mean potential temperature or surface fluxes), this framework introduces a structural complexity metric based on scale-by-scale energy decomposition. This capability provides an objective constraint for evaluating how models represent intermittent turbulence and wave-turbulence transitions. The underlying Julia package structure isolates the core Riemannian projection mathematics from the input grid data, ensuring compatibility with virtual model profiles once grid-mapping transformations are applied.
